\documentclass[11pt]{article}
\usepackage{fontspec}
\usepackage{xeCJK}
\setCJKmainfont[ItalicFont=STSongti-SC-Bold]{Songti SC}
\usepackage{xurl}
\usepackage[a4paper,margin=25mm]{geometry}
\usepackage{amsmath,amssymb,amsthm,mathtools}
\usepackage{xcolor}
\usepackage{booktabs,array,enumitem,fancyhdr,hyperref}
\hypersetup{colorlinks=true,linkcolor=blue!50!black,urlcolor=blue!50!black}
\setlength{\emergencystretch}{2em}
\setlength{\parskip}{0.3em}
\linespread{1.16}
\allowdisplaybreaks
\newtheorem{theorem}{定理}[section]
\newtheorem{proposition}[theorem]{命题}
\newtheorem{lemma}[theorem]{引理}
\theoremstyle{definition}
\newtheorem{definition}[theorem]{定义}
\theoremstyle{remark}
\newtheorem{remark}[theorem]{说明}
\newcommand{\dd}{\,\mathrm d}
\newcommand{\ii}{\mathrm i}
\newcommand{\Rea}{\operatorname{Re}}
\newcommand{\Ima}{\operatorname{Im}}
\newcommand{\Li}{\operatorname{Li}}
\newcommand{\supp}{\operatorname{supp}}
\pagestyle{fancy}
\fancyhf{}
\fancyhead[L]{\small 学习笔记}
\fancyhead[R]{\small 初稿}
\fancyfoot[C]{\thepage}
\setlength{\headheight}{15pt}
\author{}
\date{2026年10月22日初稿}
\makeatletter
\renewcommand{\maketitle}{%
  \thispagestyle{plain}%
  \begin{center}
    {\LARGE\bfseries \@title\par}
    \vspace{0.6em}
    {\small \@date\par}
  \end{center}
  \vspace{0.6em}}
\makeatother
\newcommand{\draftnotice}{\noindent{\small
本文为供逐篇审阅的学习笔记初稿，不是新增研究成果。
文中区分直接证明与引用结果；正文不依赖本文。
本轮已作公式、条件及排版初步核查，尚待作者逐句审定。}\par\medskip}

\XeTeXlinebreaklocale "zh"
\renewcommand{\proofname}{证明}
\date{2026年10月24日初稿}
\title{计算机能证明定理吗}
\begin{document}
\maketitle

\noindent{\small 本篇讨论计算在证明中的地位，以与 $\zeta$ 相关的计算为主要例子。
命题给出证明；文中数值均由区间（球）算术得到，并注明所依赖的软件。
可与学习笔记第 6 篇（数值计算的认证）对照阅读。}\par\medskip

\section{问题}
“计算机证明”至少有三种不同含义：
\begin{enumerate}
 \item 证明中的某些步骤由计算完成，例如验证有限多个情形或认证数值不等式；
 \item 整个证明在证明助手中形式化，由程序逐步检查；
 \item 数值实验提示结论，但不构成证明。
\end{enumerate}
前两种是证明，第三种不是。本篇讨论前两种在什么条件下成立，以及它们所依赖的信任基础。

\section{证书：验证比寻找容易}
一个计算结论能否被接受，关键在于能否附带一个容易检查的\emph{证书}。素性证明是典型例子。
\begin{proposition}[Lucas]\label{prop:lucas}
设整数 $n>2$，存在整数 $a$ 使 $a^{n-1}\equiv1\pmod n$，且对 $n-1$ 的每个素因子 $q$ 有
$a^{(n-1)/q}\not\equiv1\pmod n$。则 $n$ 是素数。
\end{proposition}
\begin{proof}
由第一个条件，$\gcd(a,n)=1$，$a$ 在 $(\mathbb Z/n\mathbb Z)^\times$ 中的阶 $d$ 整除 $n-1$。
若 $d<n-1$，则 $d$ 整除某个 $(n-1)/q$，与第二个条件矛盾。故 $d=n-1$。
而 $(\mathbb Z/n\mathbb Z)^\times$ 的阶为 $\varphi(n)\le n-1$，等号仅当 $n$ 为素数时成立。
\end{proof}
例：$n=2^{61}-1$。有
\[
 n-1=2\cdot3^2\cdot5^2\cdot7\cdot11\cdot13\cdot31\cdot41\cdot61\cdot151\cdot331\cdot1321,
\]
取 $a=37$，可以验证 $37^{n-1}\equiv1$，并且对上述 $12$ 个素因子 $q$ 均有 $37^{(n-1)/q}\not\equiv1\pmod n$。
于是 $n$ 是素数。整个证书由 $n-1$ 的分解、$a=37$ 以及 $13$ 次模幂运算组成；
因子分解的正确性只需一次乘法核对，各因子的素性可以同样递归地给出证书。

Pratt（1975）证明，每个素数 $n$ 都有长度为 $O(\log^2n)$ 位的此类递归证书，可在多项式时间内检查。
这里的要点是：寻找 $a$ 与分解 $n-1$ 可能很难，但检查证书是机械而简短的。
一个计算若能产生这样的证书，其结论的可信度就不再依赖于寻找过程的程序是否正确。

\Needspace{14\baselineskip}
\section{浮点计算如何成为证明：区间算术}
\subsection{包含原理}
\begin{proposition}\label{prop:inclusion}
设每个基本运算 $\circ\in\{+,-,\times,\div\}$ 及每个基本函数 $f$ 都有区间扩张 $\circ_I$、$f_I$，满足
\[
 x\in X,\ y\in Y\ \Longrightarrow\ x\circ y\in X\circ_IY,\qquad
 x\in X\ \Longrightarrow\ f(x)\in f_I(X)
\]
（输出区间端点向外舍入）。则对任一由这些运算构成的表达式 $E$，
把每个输入换成包含它的区间后按同样次序求值，所得区间包含 $E$ 的真值。
\end{proposition}
\begin{proof}
对表达式的构造作归纳：输入处成立；若两个子表达式的真值分别属于各自的区间，
由区间扩张的包含性质，复合后的真值属于复合后的区间。
\end{proof}
于是，若算得的区间整体位于 $(0,\infty)$ 中，就证明了真值为正。
这把“数值看起来是正的”转化为一个定理，其前提只是各基本运算区间扩张的正确性。

\subsection{两个例子}
以下用 Arb 库（通过 python-flint 0.9.0，160 位二进制精度）计算，并在224位复核。Arb 对其实现的每个函数给出带严格误差界的球
$[m\pm r]$，其正确性是本节结论的信任基础。
\begin{enumerate}
 \item \textbf{Li 系数 $\lambda_1>0$}。$\lambda_1=1+\tfrac{\gamma_E}2-\tfrac12\log(4\pi)$（学习笔记第 10 篇）。算得
 \[
  0.02309570896612<\lambda_1<0.02309570896613,
 \]
 故 $\lambda_1>0$ 得到证明。
 \item \textbf{一个临界线零点的存在性}。Hardy 函数 $Z(t)=e^{\ii\vartheta(t)}\zeta(\tfrac12+\ii t)$ 对实数 $t$ 取实值且连续，
 $\vartheta(t)=\Ima\log\Gamma(\tfrac14+\tfrac{\ii t}2)-\tfrac t2\log\pi$。算得
 \begin{align*}
  -0.000019941389&<Z(14.1347)<-0.000019941387,\\
  0.000059374978&<Z(14.1348)<0.000059374981.
 \end{align*}
 由介值定理，$\zeta$ 在临界线上 $14.1347<t<14.1348$ 处有零点。
\end{enumerate}
第二个例子证明了“存在临界线上的零点”。要证明 $\mathrm{RH}(T)$，还须证明在高度 $T$ 以下\emph{没有其他零点}，
这需要下一节的计数论证。
这里展示的十进制上下界比内部球包围宽，已包含中点显示时的舍入影响，
不把截短的小数配上原始球的极小半径。同目录的\texttt{认证示例.py}复核这些严格不等式，
并输出软件版本与工作精度；高精度普通浮点只可作交叉检查，不承担区间包含的证明。

\section{Turing 方法：证明“没有遗漏”}
对不是零点纵坐标的 $T$，
\[
 N(T)=\frac{\vartheta(T)}\pi+1+S(T),\qquad S(T)=\frac1\pi\arg\zeta(\tfrac12+\ii T).
\]
$\vartheta(T)$ 可以严格计算，困难在于 $S(T)$。
记 $m(t)$ 为已确认的、互不相交的变号区间中，右端点不超过 $t$ 的区间数；
每个区间至少含一个零点，故 $m(t)\le N(t)$。
定义已检测计数的代理
\[
 \widetilde S(t)=m(t)-1-\vartheta(t)/\pi,\qquad
 D(t)=N(t)-m(t)=S(t)-\widetilde S(t)\ge0.
\]
认证高度 $T$ 须位于所有隔离区间之外。此时对 $t\ge T$，
新增的已检测区间全部位于 $(T,t]$ 内，因而 $D(t)\ge D(T)$；
这并不表示 $D$ 在任意高度都单调。

由 Turing 型积分估计，对适用范围内的 $T'>T$，可引用显式上界
\[
 \int_T^{T'}S(t)\dd t\le B(T,T')
\]
（常数及高度条件见 Trudgian 2011，本书第19章给出相应实例）。
若区间计算给出 $L\le\int_T^{T'}\widetilde S(t)\dd t$，并验证
\[
 L+(T'-T)>B(T,T'),
\]
则不能有 $D(T)\ge1$：否则整个区间上 $D(t)\ge1$，
将迫使 $\int_T^{T'}S(t)\dd t\ge L+(T'-T)>B(T,T')$，矛盾。
于是 $N(T)=m(T)$，每个已检测区间恰含一个单零点，且没有其他零点。
关键不是符号变化与预期“大致一致”，而是带严格包围和适用条件的这一不等式。
Platt 与 Trudgian（2021）用这一方法和 Arb 区间运算证明了 $\mathrm{RH}(3\cdot10^{12})$。

\section{依赖大规模计算的证明}
\begin{center}
\small
\renewcommand{\arraystretch}{1.25}
\begin{tabular}{>{\raggedright\arraybackslash}p{0.2\textwidth}>{\raggedright\arraybackslash}p{0.43\textwidth}>{\raggedright\arraybackslash}p{0.27\textwidth}}
\toprule
结果 & 计算承担的部分 & 后续核查\\
\midrule
四色定理（Appel--Haken 1976） & 检查约 $1900$ 个不可避免构形的可约性 & Robertson 等 1997 重做；Gonthier 2005 在 Coq 中形式化\\
Kepler 猜想（Hales 1998/2005） & 线性规划与区间运算验证大量不等式 & Flyspeck 项目 2014 完成形式化\\
三素数定理（Helfgott 2013） & 对 $n\le8.875\times10^{30}$ 的奇数直接验证（Helfgott--Platt） & 以严格区间运算实现\\
$\mathrm{RH}(3\cdot10^{12})$（Platt--Trudgian 2021） & 零点隔离与 Turing 计数 & Arb 球算术\\
\bottomrule
\end{tabular}
\end{center}
三素数定理的结构值得注意：解析方法证明了对充分大的 $n$ 结论成立，并给出明确的“充分大”界限；
计算覆盖其余的有限范围。两部分合起来才是完整证明。
这与学习笔记第 13 篇的观点一致：一个 $\Pi_1$ 命题的有限部分可以由计算证明，无穷部分则须由论证完成。

\section{形式化证明}
证明助手（Coq/Rocq、Isabelle、HOL Light、Lean 等）要求把证明写成形式系统中的推导，
由一个很小的\emph{内核}逐步检查。可信度于是集中到三处：
\begin{enumerate}
 \item 内核的正确性。内核通常只有数千行，可以独立审查；有的系统允许导出证明项由另一程序再次检查（de Bruijn 准则）。
 \item 形式化陈述是否准确表达了原命题。这一步只能由人检查，是形式化中最常见的风险。
 \item 所依赖的公理与库定义。
\end{enumerate}
例如，Lean 的 mathlib 库已经形式化了 $\zeta$ 函数、函数方程与 RH 的陈述（Loeffler--Stoll 2025）。
这意味着 RH 的\emph{陈述}已无歧义地写成机器可检查的形式；任何声称的证明都可以在此基础上被检查。

\section{讨论}
\begin{enumerate}
 \item 计算参与的证明与传统证明的区别不在于是否“由人完成”，而在于检查方式。
 传统证明由读者逐步检查；计算证明的检查对象是算法、实现与证书。
 带证书的计算（素性证书、区间包含）与形式化证明，把需要信任的部分缩小到可以独立审查的范围。
 \item 对计算证明的常见异议是不可通览性（Tymoczko 1979）：没有人能逐步读完。
 但对长篇的人工证明，逐行通读同样不是可信度的唯一来源；独立重做与形式化对两者都适用。
 四色定理与 Kepler 猜想的后续历程说明，计算证明的可信度可以通过这些途径确立。
 \item 就 RH 而言，计算的作用有明确的边界：它可以严格证明 $\mathrm{RH}(T)$，
 可以严格证明 Li 系数、de Bruijn--Newman 常数等的有限信息，
 若 RH 不成立，原则上也可以严格证实反例。有限高度的逐段验证不能代替关于所有高度的证明，
 但这不排除计算机寻找或核查一个一般性证明；形式化证明正是后者的一种实现。
\end{enumerate}

\section*{参考}
\begingroup
\small
\begin{itemize}[itemsep=2pt,topsep=3pt,parsep=0pt]
 \item V.\ R.\ Pratt, Every prime has a succinct certificate, \emph{SIAM J.\ Comput.} 4 (1975).
 \item F.\ Johansson, Arb: efficient arbitrary-precision midpoint-radius interval arithmetic, \emph{IEEE Trans.\ Comput.} 66 (2017).
 \item A.\ M.\ Turing, Some calculations of the Riemann zeta-function, \emph{Proc.\ London Math.\ Soc.} (3) 3 (1953).
 \item T.\ Trudgian, Improvements to Turing's method, \emph{Math.\ Comp.} 80 (2011).
 \item D.\ Platt, T.\ Trudgian, The Riemann hypothesis is true up to $3\cdot10^{12}$, \emph{Bull.\ London Math.\ Soc.} 53 (2021).
 \item H.\ A.\ Helfgott, The ternary Goldbach problem, 2013/2015, \url{https://arxiv.org/abs/1501.05438}；
 H.\ A.\ Helfgott, D.\ Platt, Numerical verification of the ternary Goldbach conjecture up to $8.875\cdot10^{30}$, \emph{Exp.\ Math.} 22 (2013).
 \item G.\ Gonthier, Formal proof---the four-color theorem, \emph{Notices AMS} 55 (2008).
 \item T.\ Hales et al., A formal proof of the Kepler conjecture, \emph{Forum Math.\ Pi} 5 (2017).
 \item D.\ Loeffler, M.\ Stoll, Formalizing zeta and $L$-functions in Lean, \emph{Annals of Formalized Mathematics}, 2025. \url{https://arxiv.org/abs/2503.00959}
 \item T.\ Tymoczko, The four-color problem and its philosophical significance, \emph{J.\ Philos.} 76 (1979).
\end{itemize}
\endgroup
\end{document}
